% Round 16. Full-test activation of the audited UG/AD construction.
% ACTIVATED after independent review, including the negative-output edge case.
\section{Projective-bundle actual output in all degrees and against all actual tests}
\label{r16out:section}

The finite trace theorem removes the degree boundary in the Round 15
UG/AD construction.  A separate projection-formula argument removes
its smooth-test boundary.  The latter does not follow merely from
equality of currents: it uses the ambient representative that UG
actually constructs, before passing to currents.
The independent global-height review passed the complete mathematical
source, including its negative-output degree extension, at SHA256
\texttt{d0827bc7c4e68b1e3ce1235e6ebe5e2f789f321c6183644f94a5848a62bea4fe}.
Only this activation notice differs from that reviewed version.

\begin{lemma}[Actual target factors in the universal-tree construction]
\label{r16out:target-factor}
Use the split ambient data of Round 15 Theorem \texttt{r15ug:projection}.
Suppose
\[
 \alpha=t_{(b,w)}^\star\omega,\qquad
 \beta=t_b^\star\operatorname{pr}_*(\Gamma_s\wedge\omega),
 \tag{OUT16.1}\label{r16out:local-output}
\]
where $\omega$ has proper ambient support in the vertical coordinates.
For an arbitrary ordinary actual target form $\nu$, the same construction
applied to $\alpha\wedge\pi^*\nu$ gives exactly $\beta\wedge\nu$.
No condition on the delta codimension of $\nu$ is imposed.
\end{lemma}

\begin{proof}
This is local on the ordinary base.  Choose an actual presentation
$\nu=t_c^\star\nu_0$, with $c$ a finite tuple of genuine analytic units
and $\nu_0$ a Euclidean delta-form on a suitable open.  On a smaller
compact working neighborhood one may multiply $\nu_0$ by a smooth
Euclidean cutoff inside this open and extend it by zero.  The
resulting presenter has the same pullback there.  Append $c$ to $b$;
write $l$ for the coordinate projection to the old base coordinates
and $k$ for the coordinate projection to the $c$ coordinates.
With $\operatorname{pr}'$ the new vertical projection, put
\[
 \widetilde\omega=(l\times\mathrm{id})^\star\omega,
 \qquad \widetilde\nu=k^\star\nu_0.
\]
These pullbacks are along surjective Euclidean coordinate projections.
In particular they are legitimate for arbitrary delta-forms.  The
old kernel pulls back in the same way and still ignores all newly
added base coordinates.  Proper vertical support is unchanged by
adding base coordinates, and remains proper after multiplication
by $\operatorname{pr}'{}^\star\widetilde\nu$.

The affine delta projection formula gives the strict ambient identity
\[
 \operatorname{pr}'_*
  \bigl(\widetilde\Gamma_s\wedge\widetilde\omega
               \wedge\operatorname{pr}'{}^\star\widetilde\nu\bigr)
   =\operatorname{pr}'_*
        (\widetilde\Gamma_s\wedge\widetilde\omega)
                  \wedge\widetilde\nu
   =l^\star\operatorname{pr}_*(\Gamma_s\wedge\omega)
                  \wedge\widetilde\nu.
 \tag{OUT16.2}\label{r16out:ambient-projection}
\]
The second equality is only Cartesian product base change by the
extra Euclidean coordinates.  It is not a base-change assertion for
arbitrary analytic currents.  Genuine pullback through $(b,c)$
proves the asserted equality of ordinary actual forms.  All
intersections are Euclidean delta intersections, including any
positive-delta wall or vertex terms of $\nu_0$.
\end{proof}

\begin{lemma}[The AD excision works with actual target factors]
\label{r16out:excision}
In the Round 15 ambient atlas AD15.3, the same horizontal cutoff
$\rho$ is current-neutral for $\alpha'$ and for
$\alpha'\wedge\pi'^*\nu'$, for every actual target form $\nu'$.
It is also neutral for the pairing of $\alpha'$ with every compact
complementary actual target test.
\end{lemma}

\begin{proof}
On each discarded root or infinity tube the frozen actual
retraction-germ lemma gives $\alpha'=\pi'^*\beta_j$.  Thus
$\alpha'\wedge\pi'^*\nu'=\pi'^*(\beta_j\wedge\nu')$ on
the same tube, as an actual identity.  If $\eta$ is any
complementary actual test for either input, the paired product
on the $d$-dimensional base has degree $(d+1,d+1)$ and is zero.
For the second assertion this product is simply
$\beta_j\wedge\eta$.  For the first assertion it is
$\beta_j\wedge\nu'\wedge\eta$.
The vanishing is an actual dimension statement, not a claim that
a current can be restricted to a section.  Decomposing $1-\rho$
into summands supported in the finitely many chosen tubes and
multiplying by the remaining scalar cutoffs preserves the zero.
The uniform tube and compactness arguments in AD are independent
of the form degrees or of the test's delta codimension.
\end{proof}

\begin{theorem}[All-actual projective-line integration]
\label{r16out:line}
Let $S/K$ be smooth integral algebraic, let $U\subset S^{\mathrm{an}}$
be ordinary open, and let $\pi:U\times\mathbf P^1\to U$.  For all
$p,q\geq0$ there is a unique ordinary actual operator
\[
 I_\pi:B^{p,q}(U\times\mathbf P^1)\longrightarrow B^{p-1,q-1}(U)
 \tag{OUT16.3}\label{r16out:line-map}
\]
where a negative output degree denotes the zero space, such that,
for every compact complementary actual form $\eta$,
\[
 \int_U I_\pi\alpha\wedge\eta
       =\int_{U\times\mathbf P^1}\alpha\wedge\pi^*\eta.
 \tag{OUT16.4}\label{r16out:line-pairing}
\]
For every ordinary actual target form $\nu$ it satisfies
\[
 I_\pi(\alpha\wedge\pi^*\nu)=(I_\pi\alpha)\wedge\nu.
 \tag{OUT16.5}\label{r16out:line-projection}
\]
It is linear, commutes with ordinary restriction, $d'$, and $d''$,
and its current is the proper current pushforward.  Its support
is contained in the proper image of the input support.  All weights
and fields are native over the original $K$.
\end{theorem}

\begin{proof}
Theorem~\ref{r16fc:all-line} already constructs $I_\pi$ when
$p,q\geq1$, with every smooth compact mixed-test pairing, by UG,
AD, and the all-degree finite trace.  When $p=0$ or $q=0$, the same
AD presenter can still be constructed.  To verify this extension
of its stated scope, its nested neighborhoods, rational unit
replacement, root perturbations, and Euclidean partition do not
involve the form degree.  Its only degree calculation is neutral
excision, whose proof is exactly Lemma~\ref{r16out:excision} and
still applies whenever a complementary test exists.  The supplied
UG output is now the genuine affine push of a form of degree
$(p+s-1,q+s-1)$ along $\mathbf R^s$, hence is zero because one
output degree is negative.  Thus define $I_\pi\alpha=0$ in these
edge cases, using this same literal ambient formula.  We prove
\eqref{r16out:line-projection} before using any nonsmooth test.
Work on a small ordinary target neighborhood, choose the AD
finite splitting cover $a:U'\to U_0$ of full degree $e$, and
retain its entire inverse image.  Denote $a^*\nu$ by $\nu'$.
The form $\rho\alpha'$ has the single proper ambient presentation
supplied by AD.  Appending all target moments of $\nu'$ is allowed
by UG.  Lemma~\ref{r16out:target-factor} says that the split outputs
of the retained input and its product with $\nu'$ are $\beta'$
and $\beta'\wedge\nu'$, respectively.

Lemma~\ref{r16out:excision} shows, against the already permitted
smooth tests, that these are the outputs of the uncut inputs.
Thus the same splitting cover and cutoff compute both operators.
The all-degree finite trace projection formula gives
\[
 e^{-1}\operatorname{Tr}_a(\beta'\wedge a^*\nu)
       =\bigl(e^{-1}\operatorname{Tr}_a\beta'\bigr)\wedge\nu.
 \tag{OUT16.6}\label{r16out:finite-projection}
\]
This proves \eqref{r16out:line-projection} locally and hence globally.
It does not presuppose the desired all-actual nonfinite pairing.

Now let $\eta$ be an arbitrary compact complementary actual test.
The source input $\theta=\alpha\wedge\pi^*\eta$ is an actual top
form of compact support, because $\pi$ is proper.  Its output has
support inside $\operatorname{supp}\eta$: this follows from ordinary
restriction and uniqueness, since $\theta$ is zero over the
complement.  Choose an ordinary smooth compact scalar $\zeta$ equal
to one on a neighborhood of $\operatorname{supp}\eta$.  The already
proved smooth-test comparison for the top input $\theta$, tested
against $\zeta$, gives
\[
 \int_U I_\pi\theta
   =\int_U I_\pi\theta\,\zeta
   =\int_{U\times\mathbf P^1}\theta\,\pi^*\zeta
   =\int_{U\times\mathbf P^1}\alpha\wedge\pi^*\eta.
\]
The first equality uses the stated support containment and the last
uses the source support.  Formula~\eqref{r16out:line-projection}
identifies the left side with the left side of
\eqref{r16out:line-pairing}.  This proves the full actual test claim.
Proper current differentiation and actual-current injectivity prove
commutation with $d'$ and $d''$, since the input and derivative
outputs now all exist.  The even real relative degree introduces
no additional sign.  The remaining assertions follow from the
construction and ordinary uniqueness.
\end{proof}

\paragraph{The same refined-cycle calculation, including positive-delta tests.}
Here is a direct check of why the test upgrade respects the native
calibration.  In a split supplied presentation, append the entire
target-test tuple and every localization moment before forming
the common refined skeleton.  Let $T$ be the resulting base cycle
with its original native weights.  The rank and Gauss-point argument
of FT15.7 identifies the full refined source cycle as the marked-tree
family over $T$, with base weight times relative weight one.  UG15.4
then gives, on the boundary-excluded charts in that construction,
\[
 T_{\rm src}=\widetilde\Gamma_s\wedge\operatorname{pr}'{}^*[T].
 \tag{OUT16.7}\label{r16out:cycle}
\]
Neither this identity nor its native weights depend on whether the
ambient presenter $\eta_0$ of the target test is smooth or has
positive delta codimension.  Consequently, with the source presenter
$\widetilde\omega$, the all-delta affine projection formula reads
\[
\begin{split}
 \int T_{\rm src}\wedge\widetilde\omega
                           \wedge\operatorname{pr}'{}^\star\eta_0
 &=\int [T]\wedge
      \operatorname{pr}'_*(\widetilde\Gamma_s\wedge\widetilde\omega)
                           \wedge\eta_0.
\end{split}
 \tag{OUT16.8}\label{r16out:cycle-test}
\]
Every product here is an ambient delta intersection after genuine
moment refinement; there is no pullback through a piecewise affine
ray map.  Definition 4.10 of Mihatsch identifies the two sides with
the source and target actual top integrals on these charts.  The
same smooth scalar localizations used in UG restrict to compact
ordinary paired supports.  One does not infer that an arbitrary
ambient presenter avoids tropicalized analytic boundary merely
because its actual pullback is compact.  The form-level proof above
instead retains the already audited UG boundary treatment and
deduces the test upgrade from its ambient output and compact top
input.  Thus no new boundary-removal lemma is hidden in the display.

The horizontal excision in AD is neutral for the arbitrary actual
test by Lemma~\ref{r16out:excision}.  On descent, for this same test,
the native finite degree calculation is
\[
 \int_{U_0}e^{-1}\operatorname{Tr}_a\beta'\wedge\eta
 =e^{-1}\int_{U'}\beta'\wedge a^*\eta
 =e^{-1}\int_{U'\times\mathbf P^1}
           (a\times\mathrm{id})^*(\alpha\wedge\pi^*\eta)
 =\int_{U_0\times\mathbf P^1}\alpha\wedge\pi^*\eta.
 \tag{OUT16.9}\label{r16out:native}
\]
The integrand in the last finite degree formula is a compact actual
top form.  Compactness follows from properness of $\pi$ and $a$,
and the degree is the full $e$, including every component and every
constant-field contribution.

\begin{theorem}[All-actual integration on every algebraic projective bundle]
\label{r16out:bundle}
Let $S/K$ be smooth integral algebraic, let $N$ be a vector bundle
of rank $m+1\geq1$, let $U\subset S^{\mathrm{an}}$ be ordinary open,
and put $\pi:E_U=\mathbf P_S(N)|_U\to U$.  For $p,q\geq0$ there is
a unique ordinary actual operator
\[
 I_\pi:B^{p,q}(E_U)\longrightarrow B^{p-m,q-m}(U)
 \tag{OUT16.10}\label{r16out:bundle-map}
\]
with negative output spaces interpreted as zero, representing the
proper current pushforward and satisfying
\[
 \int_U I_\pi\alpha\wedge\eta=\int_{E_U}\alpha\wedge\pi^*\eta
 \tag{OUT16.11}\label{r16out:bundle-pairing}
\]
for every compact complementary actual test $\eta$.  It commutes
with $d',d''$, ordinary restriction, and multiplication by a pulled
back arbitrary actual target form.  Successive projective-bundle
operators compose to the proper current operator for the composite,
with the same full actual pairing.  No arbitrary proper pushforward
is asserted.
\end{theorem}

\begin{proof}
The rank-one case is the identity.  First take the product bundle
$U\times\mathbf P^m$.  The ordered-factor morphism
\[
 q_m:(\mathbf P^1)^m\longrightarrow\mathbf P^m
\]
is finite of native generic degree $m!$, as proved and independently
audited in Round 15 Lemma \texttt{r15ps:symmetric-quotient}.
Let $Q=\mathrm{id}_U\times q_m$ and let $P$ denote the successive
projective-line projections from $U\times(\mathbf P^1)^m$.
Define
\[
 I_m\alpha=(m!)^{-1}I_{P_m}\cdots I_{P_1}(Q^*\alpha).
 \tag{OUT16.12}\label{r16out:symmetric}
\]
All intermediate bases are ordinary opens in smooth integral
algebraic varieties.  If $p,q\geq m$, each input degree is at least
one until that factor is integrated.  If $p<m$ or $q<m$, one stage
is in the zero-output case of Theorem~\ref{r16out:line}, and all
subsequent outputs are zero.  Its projection formula, already
proved also for that case by the literal ambient formula, remains
valid after an arbitrary actual target factor is added.  Thus
Theorem~\ref{r16out:line} applies at every stage, including the
degree-zero convention needed for lower powers of a Chern form.

For an arbitrary compact actual base test $\eta$, its successive
pullbacks are compact by properness of the projective factors.
Repeated use of \eqref{r16out:line-pairing}, followed by the finite
native top-degree formula for $Q$, gives
\[
\begin{split}
 \int_U I_m\alpha\wedge\eta
  &=(m!)^{-1}\int_{U\times(\mathbf P^1)^m}
             Q^*(\alpha\wedge\pi_m^*\eta)\\
  &=\int_{U\times\mathbf P^m}\alpha\wedge\pi_m^*\eta.
\end{split}
 \tag{OUT16.13}\label{r16out:symmetric-pairing}
\]
This calculation includes the ramification of the symmetric quotient:
its degree is $m!$, not a pointwise count of distinct ordered roots.
The source integrand is compact actual top degree.  Thus all
intermediate finite and nonfinite formulas are used in their proved
test classes.  Differential and target-factor compatibility follow
at every stage of \eqref{r16out:symmetric}; the latter uses actual,
not merely PWS, projection formulas.  Current uniqueness makes the
choice and order of factors irrelevant.

Choose algebraic local trivializations of $N$.  The local outputs
above represent the same intrinsic proper current on every ordinary
overlap, so actual-current injectivity glues them.  An ordinary smooth
partition of unity on compact target test support yields the global
all-actual pairing: each localized test remains an actual form, and
its pullback has compact support.  The differential and projection
identities glue as actual identities.  Successive operators satisfy
the composite pairing by properness and the same repeated actual-test
argument; uniqueness proves the composition assertion.
\end{proof}

\begin{corollary}[All-degree diagonal Green outputs with all actual tests]
\label{r16out:green}
For $Y=S\times\mathbf P^1$ and the relative diagonal Green kernel
of Round 15, every actual $(r,r)$ input $\alpha$ has ordinary actual
representatives
\[
 G_\alpha\in B^{r-1,r-1}(Y),\qquad
 P_\alpha\in B^{r,r-1}(Y),\qquad
 Q_\alpha\in B^{r-1,r}(Y),
 \tag{OUT16.14}\label{r16out:green-output}
\]
with their prescribed pairings against all compact complementary
actual tests.  The same holds on projective-line bundles whenever
the stated diagonal kernel is supplied.  No closedness or symmetry
of $\alpha$ is required for the output assertion.
\end{corollary}

\begin{proof}
Work over a compact target neighborhood and use the frozen ordinary
immersion-germ identity near the diagonal.  It identifies the
source input there with the pullback of the target input.  For
$G$, its product with a complementary actual test has bidegree
$(\dim Y+1,\dim Y+1)$ on the target and is zero.  For $P$ the
second degree exceeds $\dim Y$, and for $Q$ the first degree does.
Thus the singular kernel and its first derivatives multiply the
zero paired actual factor on an ordinary punctured neighborhood of
the diagonal.  Choose a genuine ordinary scalar cutoff that removes
a smaller diagonal neighborhood.  The three retained source inputs
are actual of degrees $(r,r),(r+1,r),(r,r+1)$ on the full source.
Apply Theorem~\ref{r16out:bundle} in relative dimension one.
The overdegree argument shows that the cutoff changes none of the
prescribed all-actual pairings.  Current injectivity proves cutoff
independence and ordinary gluing.  No limiting regularization or
evaluation of an arbitrary distribution along the diagonal occurs.
\end{proof}

\begin{lemma}[Finite trace under an independent projective factor]
\label{r16out:finite-coordinate-base-change}
Let $q:Y\to X$ be finite dominant between normal integral separated
$K$-varieties.  For $b:X\times\mathbf P^a\to X$ and its source
counterpart $b_Y$, put $q_b=q\times\mathrm{id}$.  On ordinary opens
and in every actual bidegree one has
\[
 b^*\operatorname{Tr}_q\omega
       =\operatorname{Tr}_{q_b}(b_Y^*\omega).
 \tag{OUT16.15}\label{r16out:finite-coordinate}
\]
The formula also holds after ordinary restriction in the projective
factor.
\end{lemma}

\begin{proof}
Products with projective space preserve normality and integrality.
They preserve the generic degree, since their function fields add
independent transcendental variables.  If $q$ is Galois with group
$G$, so is $q_b$, with $G$ acting trivially on the new factor.
The actual orbit identity of Corollary~\ref{r16ft:orbit} gives
\[
 q_b^*\operatorname{Tr}_{q_b}(b_Y^*\omega)
  =\sum_{g\in G}g^*b_Y^*\omega
  =b_Y^*\sum_{g\in G}g^*\omega
  =q_b^*b^*\operatorname{Tr}_q\omega.
\]
Faithfulness of finite pullback proves \eqref{r16out:finite-coordinate}.
For arbitrary $q$, pass to its normal Galois closure and use the
trace normalization by the full degree of the map to $Y$; this
degree is unchanged by the independent factor.  Ordinary restriction
then gives the local assertion.  This is a consequence of the
explicit actual orbit identity, not of current base change.
\end{proof}

\begin{theorem}[Coordinate smooth base change for the actual operator]
\label{r16out:coordinate-base-change}
In Theorem~\ref{r16out:bundle}, let
$b:U\times\mathbf P^a\to U$, and let
$b_E:E_U\times\mathbf P^a\to E_U$ be its base change.  Write
$\pi_b:E_U\times\mathbf P^a\to U\times\mathbf P^a$.  Then
\[
 b^*I_\pi\alpha=I_{\pi_b}(b_E^*\alpha)
 \tag{OUT16.16}\label{r16out:coordinate-base-change-formula}
\]
as ordinary actual forms, in every permitted bidegree.  The same
holds when $b$ is any algebraic projective-bundle projection over
$U$, and after ordinary restrictions in its source.  No general
analytic or arbitrary proper current base change is asserted.
\end{theorem}

\begin{proof}
First let $\pi$ have relative dimension one and be the product
projection.  At a fixed base point choose exactly one AD construction
for $\alpha$: the finite moment charts, their nested Laurent domains,
the original-field perturbations, the finite splitting cover
$a:U'\to U_0$, the root functions, and the cutoff $\rho$.
Pull all of this data from $U_0$ to $U_0\times\mathbf P^a$.
The splitting cover is $a\times\mathrm{id}$ and has the same full
degree.  Its roots remain distinct, its difference units remain
units, the numerical domain margins remain the same, and the
pullback of $\rho$ has the same neutral excision property.
No approximation or new root splitting involving the independent
factor is performed.

The split retained input already has one supplied presentation
$t_{(b_0,w)}^\star\omega$ from AD.  Its pullback has exactly the
same ambient delta coefficient $\omega$ in the pulled-back base
units $b_0$ and the same fiber units $w$.  Proper support in the
vertical moment coordinates is unchanged.  Thus UG constructs
both upstairs outputs by the identical coefficient
\[
                 \operatorname{pr}_*(\Gamma_s\wedge\omega).
\]
For the enlarged base this coefficient is simply pulled through
the moments $b_0$ that ignore the new projective coordinates.
There is no full-rank requirement on this initially supplied base
tuple.  All test moments, including mixed units in the independent
factor, are appended before the calibrated source-cycle argument.
UG15.4 applies to precisely that refined base cycle; FT15.7
identifies its source family with weight one in the primitive
radius.  Hence the unchanged coefficient has every prescribed
pairing on the enlarged base.  Positive-delta tests are covered
by Theorem~\ref{r16out:line} and its ambient target-factor proof.
This proves that the upstairs output is the pullback of the
original upstairs output, as an ordinary actual form.

To descend, apply Lemma~\ref{r16out:finite-coordinate-base-change}:
\[
 (\deg a)^{-1}\operatorname{Tr}_{a\times\mathrm{id}}(b_{U'}^*\beta')
   =b^*\bigl((\deg a)^{-1}\operatorname{Tr}_a\beta'\bigr).
\]
This is exactly \eqref{r16out:coordinate-base-change-formula} for
the line.  All ordinary localizations retain the same full cover;
no individual analytic sheet is selected.

For a product projective space of any relative dimension, use the
same ordered-factor finite map $Q$ and commute the independent
factor past each of the line operators in
\eqref{r16out:symmetric}.  The algebraic diagram with $Q$ commutes,
so genuine actual pullback commutes with it.  The normalization
$m!$ is unchanged.  This proves the assertion for projective-space
products.  Algebraic local trivializations of $E_U$ prove it for
$\pi$.  Finally, if $b$ itself is a projective-bundle projection,
trivialize its vector bundle on algebraic opens and apply the
product statement.  Equality of ordinary actual forms glues on
the overlaps.  Restriction compatibility gives the final local
assertion.
\end{proof}

\paragraph{Precise new scope.}
The new statements are all-degree, all-actual-output and
all-actual-test theorems for algebraic projective bundles and their
composites, using the specific UG/AD construction.  They do not
identify arbitrary proper currents with actual forms, supply an
unproved analytic base-change theorem, or replace any positive
delta degree by PWS degree zero.  In particular a general
modification still requires its separate boundary/cohomological
analysis.  The original-field degree factors and all positive-delta
test intersections have been retained explicitly.
